\section*{PHỤ LỤC: MÃ NGUỒN PYTHON}
\phantomsection\addcontentsline{toc}{section}{\numberline{} PHỤ LỤC}

Dưới đây là các đoạn mã Python sử dụng thư viện \texttt{matplotlib} và \texttt{networkx} để tự động hóa việc vẽ các biểu đồ phân tích và sơ đồ mạng trong đồ án.

\lstset{style=customcode}

\subsection*{1. Script vẽ Biểu đồ Radar (Phân tích Trọng số)}
\begin{lstlisting}[language=Python, caption={File ve\_bieu\_do.py}]
import matplotlib.pyplot as plt
import numpy as np
import matplotlib

matplotlib.use('Agg')

def draw_radar_chart():
    # Define Criteria and Weights
    labels = [
        'Security', 'Availability', 'Cost', 'Performance',
        'Manageability', 'Scalability', 'Adaptability', 'Low Latency'
    ]
    # Values (Scale 1-10)
    stats = [9, 9, 8, 7, 7, 6, 5, 5]

    angles = np.linspace(0, 2*np.pi, len(labels), endpoint=False).tolist()
    stats += stats[:1]
    angles += angles[:1]

    fig, ax = plt.subplots(figsize=(8, 8), subplot_kw=dict(polar=True))
    ax.fill(angles, stats, color='#3498db', alpha=0.25)
    ax.plot(angles, stats, color='#3498db', linewidth=2)
    ax.set_yticklabels([])
    ax.set_xticks(angles[:-1])
    ax.set_xticklabels(labels, size=10, weight='bold')

    plt.title('Design Criteria Priority - Bao Loc Campus', size=15, y=1.1)
    plt.savefig('radar_chart_analysis.png', bbox_inches='tight', dpi=300)

if __name__ == "__main__":
    draw_radar_chart()
\end{lstlisting}

\newpage
\subsection*{2. Script vẽ Sơ đồ Logic (Topology)}
\begin{lstlisting}[language=Python, caption={File ve\_so\_do\_logic.py}]
import matplotlib.pyplot as plt
import networkx as nx
import matplotlib

matplotlib.use('Agg')

def draw_campus_topology():
    G = nx.Graph()
    # Define Nodes
    G.add_node("Core-SW1", pos=(4, 8), color='#e74c3c')
    G.add_node("Core-SW2", pos=(6, 8), color='#e74c3c')
    G.add_node("Firewall", pos=(5, 9), color='#8e44ad')
    G.add_node("Dist-SW-A", pos=(3, 6), color='#3498db')
    G.add_node("Dist-SW-B", pos=(7, 6), color='#3498db')
    
    # Define Edges
    edges = [
        ("Firewall", "Core-SW1"), ("Firewall", "Core-SW2"),
        ("Core-SW1", "Core-SW2"),
        ("Core-SW1", "Dist-SW-A"), ("Core-SW2", "Dist-SW-A"),
        ("Core-SW1", "Dist-SW-B"), ("Core-SW2", "Dist-SW-B"),
    ]
    G.add_edges_from(edges)
    
    # Draw
    pos = nx.get_node_attributes(G, 'pos')
    colors = [nx.get_node_attributes(G, 'color')[n] for n in G.nodes()]
    nx.draw(G, pos, node_color=colors, with_labels=True)
    plt.savefig('logical_topology.png')

if __name__ == "__main__":
    draw_campus_topology()
\end{lstlisting}

\subsection*{3. Script vẽ Biểu đồ Chi phí (Trade-off)}
\begin{lstlisting}[language=Python, caption={File ve\_bieu\_do\_chi\_phi.py}]
import matplotlib.pyplot as plt
import numpy as np

def draw_budget():
    categories = ['Core', 'Access', 'Wifi', 'FW', 'Cabling']
    option_a = [400, 450, 350, 200, 100] # High-end
    option_b = [200, 250, 200, 150, 100] # Optimized
    
    x = np.arange(len(categories))
    width = 0.35
    
    fig, ax = plt.subplots()
    rects1 = ax.bar(x - width/2, option_a, width, label='Option A')
    rects2 = ax.bar(x + width/2, option_b, width, label='Option B')
    
    ax.set_ylabel('Cost (Million VND)')
    ax.set_title('Budget Trade-off Analysis')
    ax.set_xticks(x)
    ax.set_xticklabels(categories)
    ax.legend()
    
    plt.savefig('budget_tradeoff.png')

if __name__ == "__main__":
    draw_budget()
\end{lstlisting}